\begingroup
\sloppy
\setlength{\parindent}{2em}

本文训练赛仅提交PDF，故在附录中给出能够复现三问核心模型与正式结果的完整程序。运行环境为Python 3.11，主要依赖为\texttt{numpy}、\texttt{pandas}、\texttt{scipy}、\texttt{matplotlib}和\texttt{openpyxl}。原始附件置于同一数据目录后，问题二、问题三程序分别通过\texttt{--data-dir}和\texttt{--project}指定数据目录与项目根目录。附录省略仅用于Origin制图、临时诊断和论文预检的辅助脚本。

\begin{longtable}{p{.39\textwidth}p{.51\textwidth}}
\toprule
核心程序 & 对应功能\\\midrule
\endfirsthead
\toprule 核心程序 & 对应功能\\\midrule\endhead
\path{shared/code/materials.py} & 问题一的SiC色散、载流子Drude修正及材料复折射率。\\
\path{shared/code/optics.py} & 问题一的Snell关系、Fresnel系数、双光束和Airy反射率正算。\\
\path{scripts/validate_q1_paper_a.py} & 问题一的物理极限、单位缩放和第三束光强比一致性检查。\\
\path{modules/30_q2/code/solve_q2_paper_a.py} & 问题二的分角度有界最小二乘、多初值筛选、厚度综合及正式结果输出。\\
\path{shared/code/data_io.py} & 问题三读取四个官方附件并执行确定性光谱预处理。\\
\path{modules/40_q3/code/solve_q3_paper_a.py} & 问题三的四参数Airy反演、第三束光强比和SiC双光束适用性检查。\\
\bottomrule
\end{longtable}

\lstset{
  language=Python,
  basicstyle=\ttfamily\footnotesize,
  breaklines=true,
  breakatwhitespace=false,
  columns=fullflexible,
  keepspaces=true,
  frame=single,
  numbers=left,
  numberstyle=\tiny,
  numbersep=6pt,
  keywordstyle=\color{CumcmBlue},
  commentstyle=\color{gray},
  showstringspaces=false,
  tabsize=4
}

\noindent\textbf{附录A：问题一材料色散与光学正算程序。}
\lstinputlisting[caption={问题一材料色散核心程序}]{../shared/code/materials.py}
\lstinputlisting[caption={问题一双光束与Airy光学正算核心程序}]{../shared/code/optics.py}
\lstinputlisting[caption={问题一物理一致性检验程序}]{../scripts/validate_q1_paper_a.py}

\noindent\textbf{附录B：问题二SiC双角度反演程序。}
\lstinputlisting[caption={问题二分角度有界反演核心程序}]{../modules/30_q2/code/solve_q2_paper_a.py}

\noindent\textbf{附录C：问题三硅多光束反演程序。}
\lstinputlisting[caption={官方光谱数据读取与预处理程序}]{../shared/code/data_io.py}
\lstinputlisting[caption={问题三Airy反演与SiC适用性检查核心程序}]{../modules/40_q3/code/solve_q3_paper_a.py}

\endgroup
